\section{Lower arities and finite-size information}\label{sec:boundaries}

The balanced comparison gives a one-sided quantitative bound on the
finite-size centers. In \eqref{eq:balanced-comparison-bound}, take
\(u_n=\mu_n\), \(A=4\), and \(\beta=\delta_k\), then let the larger
size tend to infinity. For every sufficiently large integer \(s\),
\begin{equation}\label{eq:one-sided-center-rate}
 \mu_s\le\alpha_k+C_k s^{-\delta_k},\qquad
 C_k=4\left(1+\frac{(3/2)^{\delta_k}}
                         {1-(2/3)^{\delta_k}}\right).
\end{equation}
Thus the center cannot exceed its limit by more than the stated error.
The proof gives no corresponding bound on \(\alpha_k-\mu_s\), and
hence no effective two-sided rate for approximating \(\alpha_k\).
Proposition~\ref{prop:concentration} controls fluctuations about
\(\EE V_n\); it does not place a two-sided window of width
\(n^{1-\delta_k}\) around \(\alpha_kn\). All constants concern a fixed
\(k\), and are not uniform when \(k\) grows with \(n\).

\begin{remark}[The two-clause boundary]\label{rem:two}
Chv\'atal and Reed~\cite[Theorems~3--4]{CR} proved that independent
uniform clauses on two distinct variables, sampled with replacement,
are satisfiable with probability tending to one below density \(1\)
and unsatisfiable with probability tending to one above density \(1\).
Goerdt~\cite{Goerdt} obtained the threshold independently.
Thus \(\alpha_2=1\) in the same proper-clause convention. This is a
separate theorem: the bound \(O_k(n^{1+2/k})\) is only \(O(n^2)\) at
\(k=2\) and gives no sublinear concentration window there.
\end{remark}

\begin{remark}[The unit-clause boundary]\label{rem:unit}
For \(k=1\), the threshold on the linear-density scale is zero.
Fix \(c>0\), choose \(0<c'<c\), and Poissonize the number of independent
uniform unit clauses at mean \(c'n\). For each variable, the positive
and negative occurrence counts are independent \(\Pois(c'/2)\)
variables, independently across variables. The formula is satisfiable
exactly when no variable receives both signs, an event of probability
\[
 \bigl(2e^{-c'/2}-e^{-c'}\bigr)^n\longrightarrow0.
\]
The Poisson clause count is at most \(\lfloor cn\rfloor\) with
probability tending to one. Monotonicity in the clause count therefore
implies that the fixed-count unit-clause formula is unsatisfiable with
probability tending to one at every \(c>0\). At density zero it is
empty and satisfiable.
\end{remark}
